% Options for packages loaded elsewhere
% Options for packages loaded elsewhere
\PassOptionsToPackage{unicode}{hyperref}
\PassOptionsToPackage{hyphens}{url}
\PassOptionsToPackage{dvipsnames,svgnames,x11names}{xcolor}
%
\documentclass[
  french,
  11pt,
  a4paper,
  twoside,
  article,
  titlepage=false]{memoir}
\usepackage{xcolor}
\usepackage{amsmath,amssymb}
\setcounter{secnumdepth}{-\maxdimen} % remove section numbering
\usepackage{iftex}
\ifPDFTeX
  \usepackage[T1]{fontenc}
  \usepackage[utf8]{inputenc}
  \usepackage{textcomp} % provide euro and other symbols
\else % if luatex or xetex
  \usepackage{unicode-math} % this also loads fontspec
  \defaultfontfeatures{Scale=MatchLowercase}
  \defaultfontfeatures[\rmfamily]{Ligatures=TeX,Scale=1}
\fi
\usepackage{lmodern}
\ifPDFTeX\else
  % xetex/luatex font selection
\fi
% Use upquote if available, for straight quotes in verbatim environments
\IfFileExists{upquote.sty}{\usepackage{upquote}}{}
\IfFileExists{microtype.sty}{% use microtype if available
  \usepackage[]{microtype}
  \UseMicrotypeSet[protrusion]{basicmath} % disable protrusion for tt fonts
}{}
\makeatletter
\@ifundefined{KOMAClassName}{% if non-KOMA class
  \IfFileExists{parskip.sty}{%
    \usepackage{parskip}
  }{% else
    \setlength{\parindent}{0pt}
    \setlength{\parskip}{6pt plus 2pt minus 1pt}}
}{% if KOMA class
  \KOMAoptions{parskip=half}}
\makeatother
% Make \paragraph and \subparagraph free-standing
\makeatletter
\ifx\paragraph\undefined\else
  \let\oldparagraph\paragraph
  \renewcommand{\paragraph}{
    \@ifstar
      \xxxParagraphStar
      \xxxParagraphNoStar
  }
  \newcommand{\xxxParagraphStar}[1]{\oldparagraph*{#1}\mbox{}}
  \newcommand{\xxxParagraphNoStar}[1]{\oldparagraph{#1}\mbox{}}
\fi
\ifx\subparagraph\undefined\else
  \let\oldsubparagraph\subparagraph
  \renewcommand{\subparagraph}{
    \@ifstar
      \xxxSubParagraphStar
      \xxxSubParagraphNoStar
  }
  \newcommand{\xxxSubParagraphStar}[1]{\oldsubparagraph*{#1}\mbox{}}
  \newcommand{\xxxSubParagraphNoStar}[1]{\oldsubparagraph{#1}\mbox{}}
\fi
\makeatother


\usepackage{longtable,booktabs,array}
\usepackage{calc} % for calculating minipage widths
% Correct order of tables after \paragraph or \subparagraph
\usepackage{etoolbox}
\makeatletter
\patchcmd\longtable{\par}{\if@noskipsec\mbox{}\fi\par}{}{}
\makeatother
% Allow footnotes in longtable head/foot
\IfFileExists{footnotehyper.sty}{\usepackage{footnotehyper}}{\usepackage{footnote}}
\makesavenoteenv{longtable}
\usepackage{graphicx}
\makeatletter
\newsavebox\pandoc@box
\newcommand*\pandocbounded[1]{% scales image to fit in text height/width
  \sbox\pandoc@box{#1}%
  \Gscale@div\@tempa{\textheight}{\dimexpr\ht\pandoc@box+\dp\pandoc@box\relax}%
  \Gscale@div\@tempb{\linewidth}{\wd\pandoc@box}%
  \ifdim\@tempb\p@<\@tempa\p@\let\@tempa\@tempb\fi% select the smaller of both
  \ifdim\@tempa\p@<\p@\scalebox{\@tempa}{\usebox\pandoc@box}%
  \else\usebox{\pandoc@box}%
  \fi%
}
% Set default figure placement to htbp
\def\fps@figure{htbp}
\makeatother



\ifLuaTeX
\usepackage[bidi=basic]{babel}
\else
\usepackage[bidi=default]{babel}
\fi
% get rid of language-specific shorthands (see #6817):
\let\LanguageShortHands\languageshorthands
\def\languageshorthands#1{}


\setlength{\emergencystretch}{3em} % prevent overfull lines

\providecommand{\tightlist}{%
  \setlength{\itemsep}{0pt}\setlength{\parskip}{0pt}}



 


\makepagestyle{style-td1}
\makeevenhead{style-td1}{\sffamily\small Intégration et Probabilités}{\sffamily B}{\sffamily\small CC I MIASHS 1 \& 2}
\makeoddhead{style-td1}{\sffamily\small Intégration et Probabilités}{\sffamily B}{\sffamily\small CC I MIASHS 1 \& 2}
\makeevenfoot{style-td1}{\sffamily\small MA1AY010}{\thepage}{\sffamily\small L3 MIASHS}
\makeoddfoot{style-td1}{\sffamily\small L3 MIASHS}{\thepage}{\sffamily\small MA1AY010}
\pagestyle{style-td1}
\makeatletter
\@ifpackageloaded{tcolorbox}{}{\usepackage[skins,breakable]{tcolorbox}}
\@ifpackageloaded{fontawesome5}{}{\usepackage{fontawesome5}}
\definecolor{quarto-callout-color}{HTML}{909090}
\definecolor{quarto-callout-note-color}{HTML}{0758E5}
\definecolor{quarto-callout-important-color}{HTML}{CC1914}
\definecolor{quarto-callout-warning-color}{HTML}{EB9113}
\definecolor{quarto-callout-tip-color}{HTML}{00A047}
\definecolor{quarto-callout-caution-color}{HTML}{FC5300}
\definecolor{quarto-callout-color-frame}{HTML}{acacac}
\definecolor{quarto-callout-note-color-frame}{HTML}{4582ec}
\definecolor{quarto-callout-important-color-frame}{HTML}{d9534f}
\definecolor{quarto-callout-warning-color-frame}{HTML}{f0ad4e}
\definecolor{quarto-callout-tip-color-frame}{HTML}{02b875}
\definecolor{quarto-callout-caution-color-frame}{HTML}{fd7e14}
\makeatother
\makeatletter
\@ifpackageloaded{caption}{}{\usepackage{caption}}
\AtBeginDocument{%
\ifdefined\contentsname
  \renewcommand*\contentsname{Table des matières}
\else
  \newcommand\contentsname{Table des matières}
\fi
\ifdefined\listfigurename
  \renewcommand*\listfigurename{Liste des Figures}
\else
  \newcommand\listfigurename{Liste des Figures}
\fi
\ifdefined\listtablename
  \renewcommand*\listtablename{Liste des Tables}
\else
  \newcommand\listtablename{Liste des Tables}
\fi
\ifdefined\figurename
  \renewcommand*\figurename{Figure}
\else
  \newcommand\figurename{Figure}
\fi
\ifdefined\tablename
  \renewcommand*\tablename{Table}
\else
  \newcommand\tablename{Table}
\fi
}
\@ifpackageloaded{float}{}{\usepackage{float}}
\floatstyle{ruled}
\@ifundefined{c@chapter}{\newfloat{codelisting}{h}{lop}}{\newfloat{codelisting}{h}{lop}[chapter]}
\floatname{codelisting}{Listing}
\newcommand*\listoflistings{\listof{codelisting}{Liste des Listings}}
\usepackage{amsthm}
\theoremstyle{definition}
\newtheorem{exercise}{Exercice}[section]
\theoremstyle{remark}
\AtBeginDocument{\renewcommand*{\proofname}{Preuve}}
\newtheorem*{remark}{Remarque}
\newtheorem*{solution}{Solution}
\newtheorem{refremark}{Remarque}[section]
\newtheorem{refsolution}{Solution}[section]
\makeatother
\makeatletter
\makeatother
\makeatletter
\@ifpackageloaded{caption}{}{\usepackage{caption}}
\@ifpackageloaded{subcaption}{}{\usepackage{subcaption}}
\makeatother
\usepackage{bookmark}
\IfFileExists{xurl.sty}{\usepackage{xurl}}{} % add URL line breaks if available
\urlstyle{same}
\hypersetup{
  pdflang={fr},
  colorlinks=true,
  linkcolor={blue},
  filecolor={Maroon},
  citecolor={Blue},
  urlcolor={Blue},
  pdfcreator={LaTeX via pandoc}}


\author{}
\date{}
\begin{document}
\frontmatter


\mainmatter
\begin{tcolorbox}[enhanced jigsaw, arc=.35mm, opacityback=0, colback=white, colframe=quarto-callout-note-color-frame, leftrule=.75mm, breakable, rightrule=.15mm, bottomrule=.15mm, toprule=.15mm, left=2mm]

\vspace{-3mm}\textbf{CC I/IV :}\vspace{3mm}

\begin{itemize}
\tightlist
\item
  Licence MIASHS
\item
  20 Octobre 2025 Durée : 1 heure 15
\item
  \textbf{Intégration et Probabilités}
\end{itemize}

\end{tcolorbox}

\begin{exercise}[]\protect\hypertarget{exr-1-a}{}\label{exr-1-a}

~

\end{exercise}

\textbf{Partie I}

Dans un étang vivent \(N\geqslant 1\) poissons numérotés \(1\) à \(N\).
On pêche le poisson \(i\) avec une probabilité \(p_i\)
(\(\sum_{i=1}^N  p_i=1\)).\\
Après avoir noté le numéro du poisson pêché, on le remet (délicatement)
à l'eau. On répète cette opération \(n\geq 1\) fois. Les opérations sont
indépendantes les unes des autres.

\begin{enumerate}
\def\labelenumi{\roman{enumi}.}
\tightlist
\item
  Proposer un univers, une tribu et une loi sur la tribu pour modéliser
  ces pêches répétées.
\item
  Quelle est la probabilité que le poisson \(i\) ne soit pas pêché du
  tout ?
\item
  Quelle est la probabilité que le poisson \(i\) soit pêché au moins
  \(k\) fois?
\item
  Quelle est la probabilité que le poisson pêché lors de l'épreuve \(n\)
  ait été pêché exactement une fois lors des \(n-1\) épreuves
  précédentes.
\end{enumerate}

\textbf{Partie II}

Dans le même étang, au lieu de répéter l'épreuve \(n\) fois, on tire
d'abord \(X\) selon une loi de Poisson de paramètre \(n\). Si \(X=k\),
on effectue ensuite \(k\) pêches indépendantes selon la description plus
haut.

\begin{enumerate}
\def\labelenumi{\alph{enumi}.}
\setcounter{enumi}{21}
\tightlist
\item
  Proposer un univers, une tribu et une loi sur la tribu pour modéliser
  ces pêches répétées un nombre aléatoire de fois.
\end{enumerate}

\begin{enumerate}
\def\labelenumi{\roman{enumi}.}
\setcounter{enumi}{5}
\tightlist
\item
  Montrer que pour chaque poisson \(i\), le nombre de captures du
  poisson \(i\) suit une loi de Poisson. Préciser le paramètre de la loi
  de Poisson. Quelle est la probabilité que le poisson \(i\) soit pêché
  exactement \(\ell\) fois?
\item
  Quelle est la probabilité que le poisson \(i\) soit pêché exactement
  \(m\) fois \emph{et} le poisson \(j\) soit pêché exactement \(m\) fois
  ?
\end{enumerate}

\emph{Rappel} : la loi de Poisson de paramètre \(\mu>0\) est définie par

\[P(\{k\}) =  \mathrm{e}^{-\mu} \frac{\mu^k}{k!} \qquad \text{ pour } k \in \mathbb{N}\, .\]

\begin{exercise}[]\protect\hypertarget{exr-3-b}{}\label{exr-3-b}

~

\end{exercise}

Répondre par \emph{vrai/faux}

\begin{itemize}
\tightlist
\item[$\square$]
  Deux événements disjoints ne peuvent pas être indépendants.
\item[$\square$]
  Un événement ne peut pas être indépendant de lui-même.
\item[$\square$]
  \((A, B, C)\) forment une famille indépendante si et seulement si
  \((A,B)\), \((A,C)\) et \((B,C)\) sont des paires d'événements
  indépendants.
\item[$\square$]
  \(\cup_n \mathbb{N}^n\) est dénombrable.
\item[$\square$]
  La tribu de parties de \(\mathbb{N}\) engendrée par les sous-ensembles
  finis de \(\mathbb{N}\) est dénombrable.
\item[$\square$]
  \(f : \mathbb{R} \to \mathbb{R}\) définie par
  \(f(x)=\mathbb{I}_{[0,1]}(x)\) est une loi de probabilité sur
  \((\mathbb{R}, \mathcal{B}(\mathbb{R}))\).
\item[$\square$]
  Une fonction de répartition atteint toujours la valeur \(1\)
\item[$\square$]
  Une fonction quantile est continue à droite en tout point.
\item[$\square$]
  L'ensemble des médianes d'une loi de probabilité sur
  \((\mathbb{R}, \mathcal{B}(\mathbb{R}))\) définit un intervalle.
\end{itemize}

\newpage{}

\begin{exercise}[]\protect\hypertarget{exr-4-b}{}\label{exr-4-b}

~

\end{exercise}

La Figure~\ref{fig-repart} représente sur l'intervalle \([-2, 7.5)\) la
fonction de répartition d'une loi de probabilité sur
\(\mathbb{R}, \mathcal{B}(\mathbb{R})\).

\begin{enumerate}
\def\labelenumi{\roman{enumi}.}
\item
  À partir de la figure et de vos connaissances, répondre par
  \emph{Vrai/Faux} aux questions suivantes.

  \begin{itemize}
  \item[$\square$]
    Cette fonction de répartition est associée à une loi discrète
  \item[$\square$]
    La loi associée possède une unique médiane
  \item[$\square$]
    La fonction quantile associée est continue sur \((0,1)\)
  \item[$\square$]
    La fonction de répartition est dérivable presque partout.
  \item[$\square$]
    La fonction de répartition est l'intégrale de sa dérivée.
  \item[$\square$]
    La loi associée est une loi sur \([0, \infty)\).
  \end{itemize}
\item
  Que vaut la fonction quantile en \(0.5\)?
\end{enumerate}

\begin{figure}

\centering{

\pandocbounded{\includegraphics[keepaspectratio]{cc1-miashs1et2-B_files/figure-pdf/fig-repart-1.pdf}}

}

\caption{\label{fig-repart}}

\end{figure}%


\backmatter


\end{document}
